\chapter{Bạn bè và các mối quan hệ}
\markboth{Bạn bè và các mối quan hệ}{Friends and Relationships}

\section{Một người chỉ hiểu bạn thật là ai sau khi mình sa cơ.}
\begin{truyen}{Một người chỉ hiểu bạn thật là ai sau khi mình sa cơ.}{Hard times reveal the true guest list.}
\chuhoa{L}{úc khá giả, điện thoại anh lúc nào cũng có người hỏi han. Khi công việc đổ vỡ, danh sách ấy rút ngắn nhanh đến mức đau lòng. Chính lúc ít người còn lại, anh mới nhận ra đâu là quan hệ, đâu là tình nghĩa.}
\textit{(When life was good, his phone was always busy. When his work collapsed, the list shrank painfully fast. In that smaller circle, he finally saw the difference between contact and true friendship.)}
\end{truyen}
\begin{baihoc}
Bạn thật thường không đông; điều muộn màng là ta cứ tưởng ai cười với mình cũng là người ở lại.
\textit{(True friends are usually few; the late lesson is realizing not everyone who smiles at you will stay.)}
\end{baihoc}
\begin{ghinhoanh}
\textbf{Short English to remember:}

Hard times reveal the true guest list.

Few is enough if it is real.
\end{ghinhoanh}
\begin{tuhocnhanh}
1. Khi em gặp khó, ai là người thật sự ở lại? \textit{(Who truly stays when you are in trouble?)}

2. Em có đang nhầm sự đông vui với sự gắn bó không? \textit{(Are you confusing popularity with real connection?)}
\end{tuhocnhanh}
\ngancach

\section{Người nói chuyện mỗi ngày chưa chắc là người hiểu mình nhất.}
\begin{truyen}{Người nói chuyện mỗi ngày chưa chắc là người hiểu mình nhất.}{Frequency is not depth.}
\chuhoa{C}{ó người nhắn tin liên tục nhưng chỉ ở bề mặt. Có người ít gặp hơn, nhưng mỗi lần gặp lại khiến ta thấy mình được hiểu thật. Nhiều người chỉ phân biệt được điều này sau khi đã lỡ đặt lòng tin sai chỗ.}
\textit{(Some people message all the time but remain on the surface. Others meet less often, yet make us feel deeply understood. Many people learn this only after placing trust in the wrong place.)}
\end{truyen}
\begin{baihoc}
Độ sâu của một mối quan hệ không đo bằng tần suất xuất hiện, mà bằng mức độ chân thành khi xuất hiện.
\textit{(The depth of a relationship is not measured by frequency, but by sincerity.)}
\end{baihoc}
\begin{ghinhoanh}
\textbf{Short English to remember:}

Frequency is not depth.

Real connection feels different.
\end{ghinhoanh}
\begin{tuhocnhanh}
1. Ai là người ít nói nhưng khiến em thấy được hiểu? \textit{(Who speaks little but makes you feel understood?)}

2. Em đang đầu tư vào mối quan hệ sâu hay mối quan hệ ồn? \textit{(Are you investing in deep relationships or noisy ones?)}
\end{tuhocnhanh}
\ngancach

\section{Một người kể hết chuyện riêng vì nghĩ thân là đủ.}
\begin{truyen}{Một người kể hết chuyện riêng vì nghĩ thân là đủ.}{Trust should grow with proof.}
\chuhoa{V}{ì muốn được chia sẻ, cô mở lòng quá nhanh. Khi chuyện riêng thành đề tài sau lưng, cô mới hiểu tin tưởng cũng cần thời gian và quan sát, không chỉ cảm giác thân lúc đầu.}
\textit{(Wanting connection, she opened up too quickly. When her private story became gossip, she learned that trust needs time and evidence, not only early warmth.)}
\end{truyen}
\begin{baihoc}
Tin người là cần, nhưng đặt lòng tin sai chỗ thường khiến ta học một bài rất đau.
\textit{(Trusting people matters, but placing trust wrongly often teaches a painful lesson.)}
\end{baihoc}
\begin{ghinhoanh}
\textbf{Short English to remember:}

Trust should grow with proof.

Warmth is not always safety.
\end{ghinhoanh}
\begin{tuhocnhanh}
1. Em thường tin người theo cảm xúc hay theo quan sát? \textit{(Do you trust people based on feeling or observation?)}

2. Ranh giới nào em cần giữ kỹ hơn với người khác? \textit{(What boundary do you need to protect more carefully?)}
\end{tuhocnhanh}
\ngancach

\section{Một người mất bạn chỉ vì cố hơn thua trong lời nói.}
\begin{truyen}{Một người mất bạn chỉ vì cố hơn thua trong lời nói.}{Being right can be too expensive.}
\chuhoa{T}{rận cãi vã bắt đầu từ chuyện nhỏ, nhưng ai cũng muốn thắng. Khi mối quan hệ gãy đi, anh mới hiểu có những lần đúng về lý mà sai về cách sống, và cái giá của việc hơn thua là một người từng rất quý.}
\textit{(The argument began with something small, but both sides wanted to win. When the friendship broke, he understood that one can be correct in logic yet wrong in living, and the cost of winning can be someone precious.)}
\end{truyen}
\begin{baihoc}
Không phải cuộc tranh cãi nào thắng cũng đáng; có những chiến thắng làm mình nghèo đi.
\textit{(Not every argument is worth winning; some victories leave us poorer.)}
\end{baihoc}
\begin{ghinhoanh}
\textbf{Short English to remember:}

Being right can be too expensive.

Choose the person over the ego.
\end{ghinhoanh}
\begin{tuhocnhanh}
1. Em có đang cố thắng ở một mối quan hệ nào không? \textit{(Are you trying to win in a relationship right now?)}

2. Điều gì quan trọng hơn: chứng minh mình đúng hay giữ nhau còn gần? \textit{(What matters more: proving you are right or staying close?)}
\end{tuhocnhanh}
\ngancach

\section{Một người luôn xuất hiện khi cần, nhưng chẳng ai nhớ lúc anh ổn.}
\begin{truyen}{Một người luôn xuất hiện khi cần, nhưng chẳng ai nhớ lúc anh ổn.}{Do not become useful only.}
\chuhoa{A}{nh quen giúp bạn bè dọn hậu quả, giải quyết rắc rối, lắng nghe khủng hoảng. Lâu dần, anh nhận ra nhiều người chỉ tìm đến khi cần dùng, còn niềm vui hay sự bình yên của anh thì chẳng ai hỏi.}
\textit{(He was always there to fix problems, clean up messes, and listen to crises. Over time, he realized many people came only when they needed something, while no one asked about his peace or joy.)}
\end{truyen}
\begin{baihoc}
Tử tế không có nghĩa là để mình bị dùng mãi; trưởng thành là biết phân biệt người cần mình và người chỉ cần lợi ích từ mình.
\textit{(Kindness does not mean allowing yourself to be used forever; maturity means knowing the difference between people who need you and people who only need your benefits.)}
\end{baihoc}
\begin{ghinhoanh}
\textbf{Short English to remember:}

Do not become useful only.

Kindness still needs boundaries.
\end{ghinhoanh}
\begin{tuhocnhanh}
1. Trong các mối quan hệ, em có đang chỉ được nhớ đến khi có việc không? \textit{(In your relationships, are you remembered only when there is a need?)}

2. Em cần đặt lại ranh giới với ai? \textit{(With whom do you need to reset your boundaries?)}
\end{tuhocnhanh}
\ngancach